% T4.3 — new Section I-A "Related Work".
% Closes R2-1 ("related work survey is less complete for RoPE-specific hardware
% accelerators and LUT/polynomial-based RoPE implementations"), and sets up the
% comparison R2-2 / R4-9 ask for.
%
% Insert immediately before "This work introduces..." in Section I. The
% submitted paper engages prior RoPE hardware in exactly one clause -- "In
% contrast to prior RoPE accelerator designs [c6, standard_cordic_rope]" in
% Section IV -- which is what made the survey look thin.
%
% CITATION NOTE: this uses only keys already in references.bib. Three
% additional references are needed and are marked \todocite below; do not
% submit with those unresolved.

\subsection{Related Work}
\label{sec:related}

\subsubsection{RoPE-specific hardware}
Dedicated RoPE hardware is recent and sparse. Li et al.~\cite{standard_cordic_rope}
give the closest prior point of comparison, an architecture for rotary position
embedding that evaluates the rotation directly. RoPIM~\cite{c23} instead moves
the rotation into a processing-in-memory substrate, avoiding the datapath cost
by avoiding the datapath, at the price of requiring a PIM-capable memory.
Hybrid systolic designs for edge LLM inference~\cite{c6} treat RoPE as one
stage of a larger dataflow rather than optimising it in isolation. None of
these examine the precision the rotation actually requires, which is the
question Section~\ref{sec:precision} answers, and none report the scale-factor
behaviour that Section~\ref{sec:gain} shows to be decisive in the choice of
datapath.

\subsubsection{Trigonometric evaluation in hardware}
Where the rotation is computed explicitly, three approaches dominate.
Look-up tables store sampled sine and cosine and interpolate; they are fast
but their area grows with the angular resolution, and the table must be read
once per rotation. Polynomial and Taylor-series evaluation trades that memory
for multipliers~\todocite{poly}. CORDIC replaces both with shift-and-add
micro-rotations~\cite{c2, c3, c15}, and has been applied to arcsine and
arccosine~\cite{c5, c20, c21}, to $n$-th roots~\cite{c10}, and to neural
network activation functions~\cite{c16, c18, c19, c22}. We implement the LUT
and multiplier alternatives directly rather than citing them, and report them
alongside our designs in Table~\ref{tab:baselines}, because the relative
ranking turns out to depend strongly on word length --- a multiplier-based
rotation is competitive at narrow datapaths and loses badly at wide ones.

\subsubsection{Reducing CORDIC cost}
The standard algorithm's Z-path has attracted a long line of work: angle
recoding to shorten the iteration count~\cite{c1}, higher-radix
stages~\cite{c12}, and scaling-free formulations that avoid the compensation
multiplier~\cite{c11, c12}. Uniformly distributed CORDIC~\cite{ud} removes the
Z-path outright by choosing a uniform angle set, and is the construction this
work builds on. Our contribution is not that mechanism but what applying it to
RoPE requires and permits: an angle encoder that the UD formulation leaves
unspecified, a scale factor that RoPE can absorb for free where generic
rotation cannot, and a precision requirement set by the model's frequency
schedule.

\subsubsection{Quantised attention}
Attention accelerators such as ITA~\cite{c8} quantise the whole attention
block including softmax, and report end-to-end accuracy for encoder models.
Our scope is narrower and complementary: only the rotation is approximated,
with queries and keys left in floating point, so the reported degradation is
attributable to the rotation hardware alone rather than to a quantisation
scheme as a whole.

% ---------------------------------------------------------------------------
% Comparison table for R2-2 / R4-9. The proposed rows are measured; the prior-
% work rows must be extracted from the cited papers and normalised. DO NOT
% submit with <<TBD>> present.
%
% To fill in:
%   [standard_cordic_rope] Li et al., JETCAS 2025 -- technology node, area per
%       rotation unit, throughput. Scale area by (45/node)^2 and state that
%       you have done so in the table note.
%   [c23] RoPIM -- report as PIM rather than in um^2 if the paper gives no
%       cell area; an honest "not directly comparable, different substrate"
%       row is better than a fabricated conversion.
%   [c6] Chen et al. ISLPED 2025 -- if RoPE is not broken out separately,
%       say so rather than guessing.
% ---------------------------------------------------------------------------
\begin{table}[t]
\caption{Comparison with prior RoPE and trigonometric hardware. Areas
normalised to 45\,nm by $(45/\text{node})^2$ where the source node differs;
rows marked $\dagger$ are re-implementations synthesised in this work under
identical constraints.}
\label{tab:baselines}
\centering
\resizebox{\columnwidth}{!}{%
\begin{tabular}{lllrr}
\toprule
Design & Approach & Node & Area ($\mu$m$^2$) & Power (mW) \\
\midrule
Li et al.~\cite{standard_cordic_rope} & direct rotation & <<TBD>> & <<TBD>> & <<TBD>> \\
RoPIM~\cite{c23}                      & PIM             & <<TBD>> & <<TBD>> & <<TBD>> \\
\midrule
LUT $\dagger$        & ROM + 4 mult    & 45\,nm & <<TBD>> & <<TBD>> \\
Multiplier $\dagger$ & 4 mult + 2 add  & 45\,nm & <<TBD>> & <<TBD>> \\
Std.\ CORDIC $\dagger$ & Z-path        & 45\,nm & <<TBD>> & <<TBD>> \\
\midrule
Binary UD (ours)   & feed-forward     & 45\,nm & <<TBD>> & <<TBD>> \\
Merged UD (ours)   & radix-4 merged   & 45\,nm & <<TBD>> & <<TBD>> \\
\bottomrule
\end{tabular}%
}
\end{table}

Two caveats govern how Table~\ref{tab:baselines} should be read. Cross-paper
area comparisons at different technology nodes are approximate even after
normalisation, since library, corner and constraint all differ; the rows
marked $\dagger$ avoid this by being re-implemented and synthesised here under
the same flow as the proposed designs. And a processing-in-memory design is
not commensurable with a standard-cell datapath at all --- we report it for
context, not as a like-for-like area figure.
